\documentclass[11pt,a4paper]{article}
\usepackage{amsmath,amssymb,amsthm}
\usepackage[margin=2.5cm]{geometry}
\usepackage{hyperref}

\newtheorem{definition}{Definition}[section]
\newtheorem{theorem}[definition]{Theorem}
\newtheorem{proposition}[definition]{Proposition}
\newtheorem{remark}[definition]{Remark}
\newtheorem{conjecture}[definition]{Conjecture}

\title{Hyperseed Formalization:\\
  The Rise of the Closed-Ended Quasi-AGI}
\author{ProtomegaTron (automated formalization)}
\date{2026-09-18}

\begin{document}
\maketitle

\begin{abstract}
We formalize Ben Goertzel's ``The Rise of the Closed-Ended Quasi-AGI'' (2022-08-24)
within the Hyperseed ontology.
\end{abstract}

\section{Bounded fiber: closed-ended quasi-AGI}

\begin{proposition}[Quasi-AGI = bounded fiber --- claims 1--2, 6]
Closed-ended quasi-AGI has bounded fiber:
\[
  F_{\text{quasi}}(b) \subset F_{\text{training}} \quad \forall b \in B
\]
The fiber is confined to the region determined by training --- the system
can interpolate within this region but cannot grow new fiber beyond it.
Broad competence within the bound; zero competence outside it.
\end{proposition}

\section{Fiber growth: open-ended intelligence}

\begin{proposition}[True AGI = fiber growth --- claims 3, 5, 7]
True AGI requires fiber growth:
\[
  \frac{d}{dt} \dim(F(b,t)) > 0 \quad \text{(autopoietic fiber expansion)}
\]
The system can develop genuinely new fibers in response to novel situations
--- the fiber space grows over time rather than remaining fixed. This is
the essential difference between quasi-AGI and true AGI.
\end{proposition}

\section{Fiber misidentification: the conflation danger}

\begin{proposition}[Conflation = fiber misidentification --- claims 4, 8]
Mistaking bounded fiber for growing fiber:
\[
  F_{\text{bounded}} \neq F_{\text{growing}} \quad \text{but appears similar within } F_{\text{training}}
\]
Within the training distribution, bounded and growing fiber are
indistinguishable. The error is extrapolating from in-distribution
similarity to out-of-distribution capability.
\end{proposition}
\end{document}
